\subsection{\texorpdfstring{\(\mathcal{L}_{\mathfrak{S}}(\mathbf{E};\mathbf{F})\) 的完备子集}{\textbackslash mathcal\{L\}\_\{\textbackslash mathfrak\{S\}\}(\textbackslash mathbf\{E\};\textbackslash mathbf\{F\}) 的完备子集}}

命题 11. --- 设 E 和 F 是两个局部凸空间，S 是由有界子集组成的 E 的一个覆盖。若 F 是 Hausdorff 且拟完备的（III, 第 8 页），则 \(\mathcal{L}(\mathrm{E};\mathrm{F})\) 中对 S-拓扑为闭的每个等度连续子集 H 都是 \(\mathcal{L}_{\mathfrak{S}}(\mathrm{E};\mathrm{F})\) 的一个完备一致子空间。

由于 H 在 \(\mathcal{L}_{\mathfrak{S}}(\mathrm{E};\mathrm{F})\) 中有界（III, 第 22 页, 命题 9），且在 \(F^{E}\) 中对 S-拓扑为闭（III, 第 16 页, 命题 4），这由 GT, X, § 1, No.~5 的推论 3 得出。

注 1. --- 设 M 是 \(\mathcal{L}_{\mathfrak{S}}(\mathrm{E};\mathrm{F})\) 的一个完备一致子空间。对 E 的有界子集所成的每个集 \(\mathfrak{S}' \supset \mathfrak{S}\)，\(\mathfrak{S}'\) -拓扑细于 \(\mathfrak{S}\) -拓扑（在 \(\mathcal{L}(\mathrm{E};\mathrm{F})\) 上）；另一方面，存在 \(\mathfrak{S}'\) -拓扑的 0 的一个基本邻域系，其中的邻域对简单收敛拓扑是闭的（III, 第 13 页, 注 2），从而更对 \(\mathfrak{S}\) -拓扑是闭的。我们断定（GT, III, § 3, No.~5, 推论 1）M 对 \(\mathfrak{S}'\) -拓扑是完备的。

推论. --- 设 E 和 F 是两个局部凸空间，H 是 \(\mathcal{L}(\mathrm{E};\mathrm{F})\) 的一个等度连续子集。若 F 是 Hausdorff 且拟完备的，且 H 上的一个滤子 \(\Phi\) 在 E 的一个完全子集 T 的所有点处简单收敛，则存在从 E 到 F 中的一个连续线性映射 u，使得 \(\Phi\) 在 E 的每个预紧子集上一致收敛于 u。

因为由命题 5（III, 第 17 页），\(\Phi\) 是 E 上预紧收敛一致结构的一个 Cauchy 滤子；由命题 11，H 的闭包 \(\overline{\mathrm{H}}\)（在 \(\mathcal{L}_{pc}(\mathrm{E};\mathrm{F})\) 中）是完备的，从而 \(\Phi\) 在 E 的每个预紧子集上一致收敛到一个映射 \(u\in \overline{\mathrm{H}}\)。

注 2. --- 设 \((u_{n})\) 是从一个 Banach 空间 E 到一个 Banach 空间 F 中的连续线性映射的一个序列；可能发生这样的情况：\((u_{n}(x))\) 在 E 的一个处处稠密的向量子空间 T 的每点处都有极限，而序列 \((u_{n})\) 在赋范空间 \(\mathcal{L}(\mathrm{E};\mathrm{F})\) 中并不有界。例如，取 E 为 R 上所有在无穷远处趋于零的连续数值函数所成的空间，带上范数 \(\|f\| = \sup_{x \in \mathbf{R}} |f(x)|\)，并设 T 为连续数值函数的子空间

它由具有紧支集者组成。从 E 到 R 中的连续线性映射 \(f \mapsto nf(n)\) 的序列对所有 \(f \in T\) 收敛于 0，但在 \(\mathcal{L}_{b}(E; \mathbf{R})\) 中不有界。同一个例子表明：在空间 \(\mathcal{L}(T; \mathbf{R})\) 中，一个序列 \((v_{n})\) 可以是简单收敛的而对有界收敛拓扑不有界。

另一方面，连续线性映射 \(f \mapsto \sum_{k=1}^{n} f(k)\) 的序列

在 \(\mathcal{L}(\mathrm{T};\mathbf{R})\) 中对简单收敛拓扑是一个 Cauchy 序列，但对该拓扑在 \(\mathcal{L}(\mathrm{T};\mathbf{R})\) 中不趋于极限。

命题 12. --- 设 E 是一个有界型局部凸空间，F 是一个完备的局部凸 Hausdorff 空间，S 是 E 的有界子集的一个族且包含每个趋于 0 的序列的像。则空间 \(\mathcal{L}_{\mathfrak{S}}(\mathrm{E};\mathrm{F})\) 是完备的。

设 \(\Phi\) 是 \(\mathcal{L}_{\mathfrak{S}}(\mathrm{E};\mathrm{F})\) 中的一个 Cauchy 滤子。则 \(\Phi\) 是简单收敛拓扑下的一个 Cauchy 滤子，故在 \(\mathrm{F}^{\mathrm{E}}\) 中收敛；此外，它的极限 \(u\) 是从 E 到 F 中的一个线性映射，且 \(\Phi\) 一致收敛于 \(u\)（在 \(\mathfrak{S}\) 的每个集上）（GT, X, § 1, No.~5, 命题 5）。由此可得：在 \(u\) 之下，一个趋于零的序列的像是一个趋于零的序列，从而 \(u\) 连续，因为 E 是有界型的（III, 第 11 页, 命题 1, (iii)）。

推论 1. --- 有界型空间的强对偶是完备的。

推论 2. --- 设 E 是一个半赋范空间，F 是一个 Banach（相应地，Fréchet）空间。空间 \(\mathcal{L}_{b}(\mathrm{E};\mathrm{F})\) 是一个 Banach（相应地，Fréchet）空间。特别地，半赋范空间的对偶是一个 Banach 空间。

